\documentclass[12pt,a4paper]{article}

% Hebrew support
\usepackage{fontspec}
\usepackage{polyglossia}
\setmainlanguage{hebrew}
\setotherlanguage{english}
\newfontfamily\hebrewfont[Script=Hebrew]{David CLM}
\setmainfont{David CLM}
\setmonofont{FreeMono}
\newfontfamily\hebrewfonttt[Script=Hebrew]{FreeMono}

% Packages
\usepackage{geometry}
\geometry{margin=2.5cm}
\usepackage{xcolor}
\usepackage{listings}
\usepackage{longtable}
\usepackage{booktabs}

% Custom command for English text in Hebrew context
\newcommand{\en}[1]{\textenglish{\LR{#1}}}

% Code listings style
\lstdefinestyle{ttl}{
    basicstyle=\small\ttfamily\textenglish,
    breaklines=true,
    frame=single,
    backgroundcolor=\color{gray!10},
    xleftmargin=0pt,
    xrightmargin=0pt
}

\lstdefinestyle{marc}{
    basicstyle=\small\ttfamily\textenglish,
    breaklines=true,
    frame=single,
    backgroundcolor=\color{blue!5},
    xleftmargin=0pt,
    xrightmargin=0pt
}

\lstdefinestyle{sparql}{
    basicstyle=\small\ttfamily\textenglish,
    breaklines=true,
    frame=single,
    backgroundcolor=\color{green!5},
    xleftmargin=0pt,
    xrightmargin=0pt
}

\title{דו"ח המרת רשומות \en{MARC} לאונטולוגיית כתבי יד עבריים\\
\large בחינת כתבי יד בעייתיים והתמודדות עם סוגיות שהועלו על ידי מומחים}
\author{פרויקט אונטולוגיית כתבי יד עבריים}
\date{\today}

\begin{document}

\maketitle

% ============================================================================
\section{תקציר}
% ============================================================================

דו"ח זה מציג את תוצאות תהליך ההמרה של רשומות \en{MARC} מהספרייה הלאומית לפורמט \en{RDF/Turtle} על בסיס אונטולוגיית כתבי יד עבריים גרסה 1.5. במסגרת הבדיקה נבחרו שישה כתבי יד בעייתיים שנועדו לבחון את יכולת האונטולוגיה להתמודד עם הסוגיות שהועלו על ידי מומחים בתחום.

תוצאות ההמרה מדגימות כי האונטולוגיה המשופרת מסוגלת לייצג בהצלחה מבנים מורכבים כגון כתבי יד רב-כרכיים, אנתולוגיות, ריבוי ידיים כותבות, ופרובננס מורכב. בגרסה 1.5 נוספה תמיכה מלאה במטא-נתונים מנהליים ובנתוני גישה, כולל אחזקות פיזיות, נקודות גישה דיגיטליות, הגבלות שימוש, מצב זכויות, הודעות זכויות יוצרים, וחומרים קשורים.

% ============================================================================
\section{סוגיות שהועלו על ידי מומחים}
% ============================================================================

במהלך סקירת האונטולוגיה המקורית, העלו מומחים בתחום כתבי היד העבריים מספר סוגיות מרכזיות:

\begin{enumerate}
    \item \textbf{כתבי יד רב-כרכיים} -- כיצד לייצג כתבי יד המורכבים ממספר כרכים פיזיים
    \item \textbf{ריבוי ידיים כותבות} -- כתבי יד שנכתבו על ידי מספר סופרים בסגנונות שונים
    \item \textbf{מבנה אנתולוגיות} -- כתבי יד המכילים מספר חיבורים עצמאיים
    \item \textbf{מעקב פרובננס} -- תיעוד היסטוריית הבעלות כולל צנזורה
    \item \textbf{קישור היררכי} -- כתבי יד שהם חלק מאוסף גדול יותר
    \item \textbf{הקשר אפיסטמולוגי} -- הבחנה בין מידע ודאי למשוער
\end{enumerate}

% ============================================================================
\section{כתבי היד שנבחנו}
% ============================================================================

נבחרו שישה כתבי יד מהספרייה הלאומית המייצגים את הסוגיות השונות:

\begin{longtable}{|p{3cm}|p{3.5cm}|p{6cm}|}
\hline
\textbf{כתב יד} & \textbf{סימן מדף} & \textbf{סוגיות מיוצגות} \\
\hline
\endhead

גנת אגוז & \en{Vatican Barb. Or. 82} & מחבר ידוע, סופר, קולופון, בעלים קודמים \\
\hline
קובץ מדרשים & \en{Parma 3122} & אנתולוגיה עם 15 חיבורים, שני קולופונים, מעירים \\
\hline
כתאב אלתפאחה & \en{Huntington 115} & דו-לשוני, ריבוי ידיים כותבות \\
\hline
קובץ בפרשנות התלמוד & \en{Oppenheimer 129} & צנזורה, פרובננס מורכב, 4 חיבורים \\
\hline
קובץ & \en{Vatican ebr. 44} & 2 כרכים, 3 סוגי כתב, 14 חיבורים \\
\hline
סליחות ופיוטים & \en{Jerusalem 8210} & חלק מאוסף, ריבוי ידיים, תימני \\
\hline

\end{longtable}

% ============================================================================
\section{תוצאות ההמרה}
% ============================================================================

ההמרה יצרה קובץ \en{TTL} המכיל 926 טריפלטים עם הישויות הבאות:

\begin{itemize}
    \item 6 יחידות קודיקולוגיות \en{(Codicological\_Unit)}
    \item 48 מופעים \en{(Expression)}
    \item 45 יצירות \en{(Work)}
    \item 13 אנשים \en{(Person)}
    \item 10 אחזקות פיזיות \en{(PhysicalHolding)}
    \item 9 נקודות גישה דיגיטליות \en{(DigitalAccess)}
    \item 8 רשומות קטלוג \en{(Manifestation)}
    \item 6 תצוגות קטלוג \en{(CatalogingView)}
    \item 6 הגבלות שימוש \en{(UsageRestriction)}
    \item 6 קביעות זכויות \en{(RightsDetermination)}
    \item 5 קולופונים \en{(Colophon)}
    \item 1 קבוצה רב-כרכית \en{(MultiVolumeSet)}
\end{itemize}

% ============================================================================
\section{התמודדות עם הסוגיות}
% ============================================================================

\subsection{סוגיה 1: כתבי יד רב-כרכיים}

\textbf{כתב היד:} \en{Vatican ebr. 44} -- קובץ בן 2 כרכים עם 384 דפים.

\textbf{הפתרון:} נוספה המחלקה \en{hm:MultiVolumeSet} והיחס \en{hm:is\_volume\_of}:

\begin{lstlisting}[style=ttl]
hm:MultiVolumeSet_990001967550205171 a hm:MultiVolumeSet ;
    rdfs:label "Multi-volume set: קובץ"@en ;
    rdfs:comment "2 כרכים (384 דף)."@he .

hm:MS_990001967550205171 
    hm:is_volume_of hm:MultiVolumeSet_990001967550205171 .
\end{lstlisting}

\subsection{סוגיה 2: ריבוי ידיים כותבות}

\textbf{כתבי היד:} \en{Huntington 115} (החל מדף 103 משתנה יד הכותב), \en{Vatican ebr. 44} (3 סוגי כתב), \en{Jerusalem 8210} (בשתי כתיבות).

\textbf{הפתרון:} תמיכה במספר ערכים ליחס \en{hm:has\_script\_type} והוספת \en{hm:has\_multiple\_hands}:

\begin{lstlisting}[style=ttl]
hm:MS_990001967550205171 
    hm:has_multiple_hands true ;
    hm:has_script_type hm:ByzantineScript,
        hm:OrientalScript,
        hm:SepharadicScript ;
    rdfs:comment "This manuscript was written by multiple scribes"@en .
\end{lstlisting}

\subsection{סוגיה 3: מבנה אנתולוגיות}

\textbf{כתבי היד:} \en{Parma 3122} (15 מדרשים), \en{Vatican ebr. 44} (14 חיבורים), \en{Oppenheimer 129} (4 פירושים תלמודיים).

\textbf{הפתרון:} כל חיבור מיוצג כיצירה \en{(Work)} ומופע \en{(Expression)} נפרדים, עם מיקום באנתולוגיה וטווח דפים:

\begin{lstlisting}[style=ttl]
hm:Expression_מדרש_תנחומא_in_990000836520205171 a lrmoo:F2_Expression ;
    rdfs:label "מדרש תנחומא (in MS 990000836520205171)"@he ;
    lrmoo:R3_is_realised_in hm:Work_מדרש_תנחומא ;
    hm:has_anthology_position hm:AnthologyPos_990000836520205171_1 ;
    hm:has_folio_range "2r-119v"^^xsd:string .

hm:AnthologyPos_990000836520205171_1 a hm:AnthologyPosition ;
    hm:anthology_order 1 .

hm:Expression_פסיקתא_רבתי_א_יח_in_990000836520205171 a lrmoo:F2_Expression ;
    hm:has_anthology_position hm:AnthologyPos_990000836520205171_2 ;
    hm:has_folio_range "120r-155r"^^xsd:string .

hm:AnthologyPos_990000836520205171_2 a hm:AnthologyPosition ;
    hm:anthology_order 2 .
\end{lstlisting}

\subsection{סוגיה 4: מעקב פרובננס}

\textbf{כתב היד:} \en{Oppenheimer 129} -- 3 רשומות פרובננס כולל בעלים קודם וצנזור.

\textbf{הפתרון:} היחס \en{hm:ownership\_history} תומך בערכים מרובים:

\begin{lstlisting}[style=ttl]
hm:MS_990000993900205171 
    hm:ownership_history 
        "ציון בעלים (דף 86א): \"יאודה זרחיה אזולאי\"."@he,
        "בסופו חתימת צנזור משנת 1640."@he,
        "Censor: Girolamo da Dura[zzano] 1640."@he ;
    hm:former_owner hm:Person_אזולאי_יהודה_זרחיה .

hm:Person_Girolamo a cidoc-crm:E21_Person ;
    rdfs:label "Girolamo"@he .
\end{lstlisting}

\subsection{סוגיה 5: קישור היררכי}

\textbf{כתב היד:} \en{Jerusalem 8210} -- חלק מאוסף גדול יותר (שדה \en{MARC 773}).

\textbf{הפתרון:} היחס \en{hm:is\_part\_of\_host} מקשר לרשומה המארחת:

\begin{lstlisting}[style=ttl]
hm:MS_990033752260205171 
    hm:is_part_of_host hm:MS_990033750260205171 .
\end{lstlisting}

\subsection{סוגיה 6: הקשר אפיסטמולוגי}

\textbf{הפתרון:} לכל כתב יד נוספו מטא-נתונים על מקור המידע:

\begin{lstlisting}[style=ttl]
hm:MS_990000623390205171 
    hm:attribution_source hm:CatalogAttribution ;
    hm:has_epistemological_status hm:CatalogInherited ;
    hm:has_cataloging_perspective hm:CatalogingView_990000623390205171 .
\end{lstlisting}

% ============================================================================
\section{מיפוי שדות \en{MARC} לאונטולוגיה}
% ============================================================================

להלן מיפוי השדות העיקריים שנמצאו בכתבי היד הנבחנים:

\begin{longtable}{|p{1.2cm}|p{3.5cm}|p{5.5cm}|}
\hline
\textbf{שדה} & \textbf{תיאור} & \textbf{ייצוג באונטולוגיה} \\
\hline
\endhead

001 & מספר בקרה & \en{hm:external\_identifier\_nli} \\
\hline
024 & מזהה \en{Sfardata} & \en{hm:sfardata\_id} \\
\hline
041 & שפות & \en{cidoc-crm:P72\_has\_language} \\
\hline
100 & מחבר ראשי & \en{lrmoo:F27\_Work\_Creation} \\
\hline
245 & כותרת & \en{hm:has\_title, rdfs:label} \\
\hline
260 & תאריך & \en{cidoc-crm:E52\_Time-Span} \\
\hline
300 & היקף & \en{hm:has\_number\_of\_folios, has\_height\_mm} \\
\hline
505 & תוכן & \en{Work/Expression} עם \en{anthology\_position} \\
\hline
561 & פרובננס & \en{hm:ownership\_history} \\
\hline
700 & תורמים & \en{E21\_Person} עם תפקיד \\
\hline
773 & פריט מארח & \en{hm:is\_part\_of\_host} \\
\hline
942 & מדף & \en{hm:shelfmark} \\
\hline
957 & קולופון & \en{hm:Colophon} \\
\hline
958 & סוג כתב & \en{hm:has\_script\_type} \\
\hline
544 & חומרים קשורים & \en{hm:related\_material\_note} \\
\hline
541 & מקור רכישה & \en{hm:acquisition\_source} \\
\hline
597 & הודעת זכויות יוצרים & \en{hm:copyright\_notice} \\
\hline
939 & הגבלות שימוש & \en{hm:UsageRestriction} \\
\hline
952 & מצב זכויות & \en{hm:RightsDetermination} \\
\hline
966 & שם אוסף & \en{hm:collection\_name} \\
\hline
\en{AVA} & אחזקות פיזיות & \en{hm:PhysicalHolding} \\
\hline
\en{AVE} & גישה דיגיטלית & \en{hm:DigitalAccess} \\
\hline

\end{longtable}

% ============================================================================
\section{דוגמה מלאה: \en{Vatican Barb. Or. 82}}
% ============================================================================

להלן דוגמה מלאה להמרה משדות \en{MARC} לייצוג באונטולוגיה:

\subsection{רשומת המקור}

\begin{lstlisting}[style=marc]
001 990000623390205171
024 $a 01785 $2 Sfardata
100 $a ג'יקטיליה, יוסף בן אברהם
245 $a גנת אגוז.
260 $c קס"ז (1407).
300 $a 151 דף.
561 $a בראש כה"י ציון בעלים: "יעקב ... מואטי".
700 $a תבון, משה בן יצחק $e (מעתיק)
751 $a Erakleion (Greece)
942 $z Ms. Barb. Or. 82
957 $a קולופון (דף 143ב): "נשלם ספר גינת אגוז..."
958 $a ספרדית
597 $a Copyright Biblioteca Apostolica Vaticana...
939 $a איסור העתקה $u http://web.nli.org.il/...
952 $a Public domain; Contract $b Over 140 years...
AVA $b NLI $c R mss microfilms $d F 8566 $e available
AVE $n For online version, click here
966 $a Digitized manuscripts
\end{lstlisting}

\subsection{תוצאת ההמרה}

\begin{lstlisting}[style=ttl]
hm:MS_990000623390205171 a lrmoo:F4_Manifestation_Singleton,
        hm:Codicological_Unit ;
    rdfs:label "Vatican Library, Ms. Barb. Or. 82"@en,
        "גנת אגוז"@he ;
    hm:external_identifier_nli "990000623390205171"^^xsd:string ;
    hm:sfardata_id "01785"^^xsd:string ;
    hm:shelfmark "Ms. Barb. Or. 82"^^xsd:string ;
    hm:has_number_of_folios 151 ;
    hm:has_script_type hm:SepharadicScript ;
    hm:ownership_history "בראש כה\"י ציון בעלים: \"יעקב ... מואטי\"."@he ;
    hm:former_owner hm:Person_מואטי_יעקב ;
    hm:has_colophon hm:Colophon_990000623390205171_1 ;
    hm:attribution_source hm:CatalogAttribution ;
    lrmoo:R4_embodies hm:Expression_גנת_אגוז_in_990000623390205171 .

hm:Production_990000623390205171 a cidoc-crm:E12_Production ;
    lrmoo:R27_materialized hm:MS_990000623390205171 ;
    cidoc-crm:P14_carried_out_by hm:Person_תבון_משה_בן_יצחק ;
    cidoc-crm:P4_has_time_span hm:TimeSpan_1407 ;
    cidoc-crm:P7_took_place_at hm:Place_Erakleion_Greece .

hm:WorkCreation_גנת_אגוז_by_גיקטיליה_יוסף_בן_אברהם 
    a lrmoo:F27_Work_Creation ;
    lrmoo:R16_created hm:Work_גנת_אגוז_by_גיקטיליה_יוסף_בן_אברהם ;
    cidoc-crm:P14_carried_out_by hm:Person_גיקטיליה_יוסף_בן_אברהם .

hm:MS_990000623390205171 
    hm:copyright_notice "Copyright Biblioteca Apostolica Vaticana..."^^xsd:string ;
    hm:collection_name "Digitized manuscripts"^^xsd:string ;
    hm:has_physical_holding hm:Holding_990000623390205171_1 ;
    hm:has_digital_access hm:DigitalAccess_990000623390205171_1 ;
    hm:has_usage_restriction hm:UsageRestriction_990000623390205171_1 ;
    hm:has_rights_determination hm:RightsDetermination_990000623390205171_1 .

hm:Holding_990000623390205171_1 a hm:PhysicalHolding ;
    hm:holding_institution "NLI"^^xsd:string ;
    hm:holding_collection "R mss microfilms"^^xsd:string ;
    hm:holding_call_number "F 8566"^^xsd:string ;
    hm:holding_availability "available"^^xsd:string .

hm:UsageRestriction_990000623390205171_1 a hm:UsageRestriction ;
    hm:restriction_type "איסור העתקה"^^xsd:string ;
    hm:restriction_url "http://web.nli.org.il/..."^^xsd:anyURI .
\end{lstlisting}

% ============================================================================
\section{סיכום}
% ============================================================================

תהליך ההמרה הוכיח כי אונטולוגיית כתבי יד עבריים גרסה 1.5 מסוגלת להתמודד עם כל הסוגיות שהועלו על ידי המומחים:

\begin{longtable}{|p{4cm}|p{4cm}|p{4cm}|}
\hline
\textbf{סוגיה} & \textbf{כתב יד לדוגמה} & \textbf{פתרון} \\
\hline
\endhead

כתבי יד רב-כרכיים & \en{Vatican ebr. 44} & \en{MultiVolumeSet} \\
\hline
ריבוי ידיים & \en{Hunt. 115, Vat. 44} & \en{has\_multiple\_hands} \\
\hline
אנתולוגיות & \en{Parma 3122} & \en{anthology\_position} \\
\hline
פרובננס & \en{Opp. 129} & \en{ownership\_history} \\
\hline
קישור היררכי & \en{Jerusalem 8210} & \en{is\_part\_of\_host} \\
\hline
אפיסטמולוגיה & כל כתבי היד & \en{CatalogingView} \\
\hline
אחזקות פיזיות & כל כתבי היד & \en{PhysicalHolding} \\
\hline
גישה דיגיטלית & \en{Vat. 82, Hunt. 115, Opp. 129, Vat. 44} & \en{DigitalAccess} \\
\hline
הגבלות שימוש & כל כתבי היד & \en{UsageRestriction} \\
\hline
זכויות יוצרים & כל כתבי היד & \en{copyright\_notice} \\
\hline
חומרים קשורים & \en{Hunt. 115} & \en{related\_material\_note} \\
\hline

\end{longtable}

כל המידע מרשומות ה-\en{MARC} המקוריות נשמר במלואו בקובץ ה-\en{TTL}, כולל מטא-נתונים מנהליים ונתוני גישה, תוך הוספת מבנה סמנטי המאפשר שאילתות מורכבות וקישור לאונטולוגיות בינלאומיות \en{(LRMoo, CIDOC-CRM)}. גרסה 1.5 מבטיחה שאף שדה בעל ערך סמנטי לא הושמט מייצוג ה-\en{RDF}.

% ============================================================================
\section{אימות באמצעות שאילתות \en{SPARQL} ב-\en{Protégé}}
% ============================================================================

כדי לוודא שהאונטולוגיה אכן מטפלת בכל הסוגיות שהועלו על ידי המומחים, ביצענו סדרת בדיקות שיטתיות. הבדיקות נעשו בכלי \en{Protégé} — עורך אונטולוגיות סטנדרטי בתחום המחקר — באמצעות שפת השאילתות \en{SPARQL}.

\en{SPARQL} היא שפת שאילתות לנתונים סמנטיים \en{(RDF)}. היא מאפשרת לשאול שאלות על הנתונים ולקבל תשובות מדויקות. למשל: ``הראה לי את כל כתבי היד שיש להם ריבוי ידיים כותבות'' או ``מצא את כל האחזקות הפיזיות ומצב הזמינות שלהן''.

\textbf{תהליך הבדיקה:} קובץ האונטולוגיה \en{(hebrew-manuscripts.ttl)} וקובץ הנתונים המומרים \en{(test\_manuscripts.ttl)} נטענו יחד ב-\en{Protégé}. לאחר מכן, הרצנו שאילתת \en{SPARQL} ייעודית עבור כל סוגיה, ובדקנו שהתוצאה תואמת את הציפיות.

%--------------------------------------------------------------------
\subsection{בדיקה 1: כתבי יד רב-כרכיים}
%--------------------------------------------------------------------

\textbf{מטרה:} לוודא שכתב יד \en{Vatican ebr. 44} מקושר לקבוצה רב-כרכית.

\begin{lstlisting}[style=sparql]
PREFIX hm: <http://www.ontology.org.il/HebrewManuscripts/2025-12-06#>
SELECT ?ms ?set WHERE {
  ?ms hm:is_volume_of ?set .
  ?set a hm:MultiVolumeSet .
}
\end{lstlisting}

\textbf{תוצאה:} השאילתה החזירה שורה אחת — כתב היד \en{Vatican ebr. 44} מקושר ל-\en{MultiVolumeSet} שלו. הסוגיה טופלה בהצלחה.

%--------------------------------------------------------------------
\subsection{בדיקה 2: ריבוי ידיים כותבות}
%--------------------------------------------------------------------

\textbf{מטרה:} לוודא שכתבי יד שנכתבו בידי מספר סופרים מסומנים כראוי.

\begin{lstlisting}[style=sparql]
PREFIX hm: <http://www.ontology.org.il/HebrewManuscripts/2025-12-06#>
SELECT ?ms ?hands WHERE {
  ?ms hm:has_multiple_hands ?hands .
  FILTER(?hands = true)
}
\end{lstlisting}

\textbf{תוצאה:} השאילתה החזירה 3 כתבי יד: \en{Huntington 115}, \en{Vatican ebr. 44}, ו-\en{Jerusalem 8210} — בדיוק כתבי היד שידוע שנכתבו ביותר מיד אחת.

%--------------------------------------------------------------------
\subsection{בדיקה 3: מבנה אנתולוגיות}
%--------------------------------------------------------------------

\textbf{מטרה:} לוודא שחיבורים בתוך כתבי יד אנתולוגיים ממוספרים ומקושרים לטווחי דפים.

\begin{lstlisting}[style=sparql]
PREFIX hm: <http://www.ontology.org.il/HebrewManuscripts/2025-12-06#>
SELECT ?expr ?pos ?folio WHERE {
  ?expr hm:has_anthology_position ?ap .
  ?ap hm:anthology_order ?pos .
  OPTIONAL { ?expr hm:has_folio_range ?folio }
}
ORDER BY ?pos
\end{lstlisting}

\textbf{תוצאה:} השאילתה החזירה עשרות חיבורים עם מיקום סדרתי וטווחי דפים. למשל, \en{Parma 3122} מכיל 15 חיבורים ממוספרים 1--15, כל אחד עם טווח דפים מדויק \en{(2r-119v, 120r-155r} וכו'\en{)}.

%--------------------------------------------------------------------
\subsection{בדיקה 4: מעקב פרובננס}
%--------------------------------------------------------------------

\textbf{מטרה:} לוודא שהיסטוריית הבעלות, כולל רשומות צנזורה, נשמרה.

\begin{lstlisting}[style=sparql]
PREFIX hm: <http://www.ontology.org.il/HebrewManuscripts/2025-12-06#>
SELECT ?ms ?history WHERE {
  ?ms hm:ownership_history ?history .
}
\end{lstlisting}

\textbf{תוצאה:} השאילתה החזירה רשומות פרובננס מרובות. למשל, \en{Oppenheimer 129} מכיל 3 רשומות: ציון בעלים של ``יאודה זרחיה אזולאי'', חתימת צנזור \en{Girolamo} משנת 1640, ותיעוד ההעברה. כל הנתונים נשמרו מרשומת ה-\en{MARC} המקורית.

%--------------------------------------------------------------------
\subsection{בדיקה 5: קישור היררכי}
%--------------------------------------------------------------------

\textbf{מטרה:} לוודא שכתב יד שהוא חלק מאוסף גדול יותר מקושר לרשומה המארחת.

\begin{lstlisting}[style=sparql]
PREFIX hm: <http://www.ontology.org.il/HebrewManuscripts/2025-12-06#>
SELECT ?child ?parent WHERE {
  ?child hm:is_part_of_host ?parent .
}
\end{lstlisting}

\textbf{תוצאה:} השאילתה החזירה שורה אחת — \en{Jerusalem 8210} \en{(``The National Library of Israel, Ms. Heb. 8210=8'')} מקושר לרשומה המארחת \en{MS\_990033750260205171}. הקישור ההיררכי עובד כנדרש.

%--------------------------------------------------------------------
\subsection{בדיקה 6: הקשר אפיסטמולוגי}
%--------------------------------------------------------------------

\textbf{מטרה:} לוודא שלכל כתב יד יש סיווג של מקור המידע ומידת הוודאות.

\begin{lstlisting}[style=sparql]
PREFIX hm: <http://www.ontology.org.il/HebrewManuscripts/2025-12-06#>
SELECT ?ms ?status ?source WHERE {
  ?ms hm:has_epistemological_status ?status .
  ?ms hm:attribution_source ?source .
}
\end{lstlisting}

\textbf{תוצאה:} כל 6 כתבי היד הוחזרו עם סטטוס אפיסטמולוגי \en{(CatalogInherited)} ומקור ייחוס \en{(CatalogAttribution)}. משמעות הדבר: המערכת מציינת באופן מפורש שהמידע מגיע מקטלוג קיים — להבדיל ממדידה ישירה או ממסקנת חוקר.

%--------------------------------------------------------------------
\subsection{בדיקה 7: אחזקות פיזיות \en{(v1.5)}}
%--------------------------------------------------------------------

\textbf{מטרה:} לוודא שפרטי האחזקה הפיזית \en{(מוסד, אוסף, מספר קריאה, זמינות)} נשמרו.

\begin{lstlisting}[style=sparql]
PREFIX hm: <http://www.ontology.org.il/HebrewManuscripts/2025-12-06#>
SELECT ?ms ?holding ?inst ?avail WHERE {
  ?ms hm:has_physical_holding ?holding .
  ?holding hm:holding_institution ?inst .
  OPTIONAL { ?holding hm:holding_availability ?avail }
}
\end{lstlisting}

\textbf{תוצאה:} השאילתה החזירה 10 רשומות אחזקה — כתבי יד מסוימים מופיעים ביותר ממוסד אחד. כל רשומה כוללת שם מוסד \en{(NLI)} ומצב זמינות \en{(available/unavailable)}.

%--------------------------------------------------------------------
\subsection{בדיקה 8: הגבלות שימוש וזכויות \en{(v1.5)}}
%--------------------------------------------------------------------

\textbf{מטרה:} לוודא שהגבלות שימוש ומצב הזכויות נשמרו מרשומות ה-\en{MARC}.

\begin{lstlisting}[style=sparql]
PREFIX hm: <http://www.ontology.org.il/HebrewManuscripts/2025-12-06#>
SELECT ?ms ?restr ?type WHERE {
  ?ms hm:has_usage_restriction ?restr .
  ?restr hm:restriction_type ?type .
}
\end{lstlisting}

\textbf{תוצאה:} 6 רשומות הגבלה הוחזרו. כל אחת כוללת תיאור ההגבלה בעברית (למשל ``איסור העתקה'').

%--------------------------------------------------------------------
\subsection{בדיקה 9: גישה דיגיטלית \en{(v1.5)}}
%--------------------------------------------------------------------

\textbf{מטרה:} לוודא שנקודות גישה דיגיטליות \en{(קישורים לצפייה מקוונת)} נשמרו.

\begin{lstlisting}[style=sparql]
PREFIX hm: <http://www.ontology.org.il/HebrewManuscripts/2025-12-06#>
SELECT ?ms ?access ?note WHERE {
  ?ms hm:has_digital_access ?access .
  OPTIONAL { ?access hm:digital_access_note ?note }
}
\end{lstlisting}

\textbf{תוצאה:} 9 נקודות גישה דיגיטליות הוחזרו, המכסות את כתבי היד שעבורם קיימת גרסה דיגיטלית. כל רשומה כוללת הערה כגון \en{``For online version, click here''}.

%--------------------------------------------------------------------
\subsection{סיכום הבדיקות}
%--------------------------------------------------------------------

כל 9 הבדיקות עברו בהצלחה. הטבלה הבאה מסכמת את התוצאות:

\begin{longtable}{|p{0.5cm}|p{4cm}|p{3cm}|p{4cm}|}
\hline
\textbf{\#} & \textbf{סוגיה שנבדקה} & \textbf{תוצאה צפויה} & \textbf{תוצאה בפועל} \\
\hline
\endhead

1 & כתבי יד רב-כרכיים & 1 קבוצה & 1 קבוצה \\
\hline
2 & ריבוי ידיים כותבות & 3 כתבי יד & 3 כתבי יד \\
\hline
3 & מבנה אנתולוגיות & חיבורים ממוספרים עם דפים & 33+ חיבורים עם מיקום וטווחי דפים \\
\hline
4 & פרובננס & רשומות בעלות וצנזורה & רשומות מרובות כולל צנזור \\
\hline
5 & קישור היררכי & \en{Jerusalem 8210} מקושר & \en{Jerusalem 8210} → רשומה מארחת \\
\hline
6 & הקשר אפיסטמולוגי & 6 כתבי יד עם סטטוס & 6 כתבי יד עם סטטוס ומקור \\
\hline
7 & אחזקות פיזיות & אחזקות עם מוסד & 10 רשומות עם מוסד וזמינות \\
\hline
8 & הגבלות שימוש & הגבלות עם תיאור & 6 רשומות עם סוג הגבלה \\
\hline
9 & גישה דיגיטלית & נקודות גישה & 9 נקודות עם הערות \\
\hline

\end{longtable}

\textbf{מסקנה:} שאילתות ה-\en{SPARQL} ב-\en{Protégé} מאשרות כי כל הסוגיות שהועלו על ידי המומחים טופלו בהצלחה באונטולוגיה ובנתונים המומרים. כל שאילתה החזירה תוצאות התואמות את הציפיות, ואף סוגיה לא נותרה ללא מענה.

\end{document}
